\documentclass[11pt,a4paper]{article}

\usepackage{amsmath,amssymb}
\usepackage{graphicx}
\usepackage{booktabs}
\usepackage{hyperref}
\usepackage{xcolor}
\usepackage{tcolorbox}
\usepackage[margin=25mm]{geometry}
\usepackage{xeCJK}
\setCJKmainfont{Hiragino Mincho ProN}
\setCJKsansfont{Hiragino Sans}

\title{素数で重力を消す？\\
\large ― 算術スペクトル変調の高校生ガイド}
\author{Wright Brothers Research Program}
\date{2026年4月}

\begin{document}
\maketitle

\begin{tcolorbox}[colback=yellow!5,colframe=orange!60!black,title=このガイドについて]
このガイドは高校生に向けて，論文
``Gravitational Coupling Control via Arithmetic Spectral Modulation''
の内容をやさしく解説します．

素数（2, 3, 5, 7, \ldots）を使って\textbf{重力を消したり，逆向きにしたり}
できるという驚くべき2つの定理を紹介します．

必要な数学は高校レベル（対数，無限級数の概念）で十分です．
微積分は「こういうものがある」程度の理解でOKです．
\end{tcolorbox}

\tableofcontents

\newpage

%% ============================================================
\section{イントロ: 重力って変えられるの？}
%% ============================================================

重力は，あなたが今この瞬間も感じている力です．
リンゴが木から落ちるのも，月が地球のまわりを回るのも，
すべて重力のしわざです．
ニュートン以来，重力の強さは「万有引力定数 $G$」という
一つの数で表されてきました．
この数は宇宙のどこでも同じ値で，変えることはできない
--- そう考えられてきました．

ところが，21世紀の数学と物理学の交差点で，
驚くべき可能性が見えてきました．
重力定数 $G$ は，実は\textbf{素数}（2, 3, 5, 7, \ldots）の
寄与の「合計」として表現できる，というのです．
もし個々の素数の寄与を操作できれば，
$G$ をゼロにしたり，マイナスにしたりできるかもしれない．
この論文は，その操作の数学的な証明を与えます．

%% ============================================================
\section{ゼータ関数とオイラー積: 素数が宇宙を作っている}
%% ============================================================

\subsection{リーマンのゼータ関数}

数学で最も有名な関数の一つに，\textbf{リーマンのゼータ関数}があります:
\begin{equation}
  \zeta(s) = \frac{1}{1^s} + \frac{1}{2^s} + \frac{1}{3^s}
  + \frac{1}{4^s} + \cdots = \sum_{n=1}^{\infty} \frac{1}{n^s}.
\end{equation}
すべての正の整数の $s$ 乗の逆数を足し合わせたものです．
たとえば $s = 2$ のとき:
\[
  \zeta(2) = 1 + \frac{1}{4} + \frac{1}{9} + \frac{1}{16} + \cdots
  = \frac{\pi^2}{6} \approx 1.645.
\]

\subsection{オイラー積: 素数だけで書き直す}

18世紀の大数学者オイラーは，この無限和を
\textbf{素数だけの積（かけ算）}に書き直せることを発見しました:
\begin{equation}
  \zeta(s) = \prod_{p:\text{素数}} \frac{1}{1 - p^{-s}}
  = \frac{1}{1 - 2^{-s}} \times \frac{1}{1 - 3^{-s}}
  \times \frac{1}{1 - 5^{-s}} \times \cdots
\end{equation}

\begin{tcolorbox}[colback=blue!5,colframe=blue!50!black,title=ポイント]
ゼータ関数は「すべての整数の和」だが，
「すべての素数の積」とも書ける．
つまり，\textbf{素数一つ一つが独立したチャンネル}として
ゼータ関数に寄与している．
各素数の寄与 $\frac{1}{1-p^{-s}}$ を
\textbf{オイラー因子}と呼ぶ．
\end{tcolorbox}

これがなぜ重要かというと，
もしゼータ関数が「真空のエネルギー」を記述しているなら，
素数ごとに真空を分解して，個別に操作できる可能性がある
--- ということです．

%% ============================================================
\section{Bost-Connes系: 素数の温度計}
%% ============================================================

1995年，数学者のBostとConnesは，
素数の世界を「熱力学」（温度・エネルギーの物理学）で
記述するシステムを構築しました．
これが\textbf{Bost-Connes系}です．

このシステムの「分配関数」（温度 $\beta$ での
全エネルギー状態の重み付き合計）は:
\begin{equation}
  Z(\beta) = \zeta(\beta).
\end{equation}
なんと，\textbf{リーマンのゼータ関数そのもの}です！

\begin{tcolorbox}[colback=green!5,colframe=green!50!black,title=温度計のたとえ]
Bost-Connes系は「素数の温度計」のようなものです．
\begin{itemize}
  \item 温度パラメータ $\beta$ が大きい（低温）:
    素数たちは秩序立って並んでいる（対称性が破れている）
  \item $\beta = 1$ で\textbf{相転移}が起きる:
    素数の秩序が崩れる臨界点
  \item $\beta$ が小さい（高温）:
    素数たちはバラバラに振る舞う（対称性が回復）
\end{itemize}
この相転移は，ゼータ関数が $s=1$ で発散する（無限大になる）
ことに対応しています．
\end{tcolorbox}

重要なのは，Bost-Connes系のエネルギー固有値が
$\lambda_n = \log n$（$n$ の対数）であることです．
このとき，「ヒートカーネル」（熱の伝わり方を記述する関数）は:
\begin{equation}
  K(t) = \sum_{n=1}^{\infty} e^{-t \cdot \log n}
  = \sum_{n=1}^{\infty} n^{-t} = \zeta(t).
\end{equation}
ヒートカーネルもゼータ関数になる．
そして\textbf{重力定数 $G$ は，このヒートカーネルの
$t=0$ 付近での振る舞いから決まる}のです．

%% ============================================================
\section{定理1: 重力ゼロ --- 2つの素数をミュートすると重力が消える}
%% ============================================================

\subsection{素数をミュートするとは？}

オイラー積の中から，特定の素数の因子を取り除くことを考えます．
たとえば素数 $2$ と $3$ を ``ミュート''（消音）すると:
\begin{equation}
  K_{\neg\{2,3\}}(t) = (1 - 2^{-t})(1 - 3^{-t}) \times \zeta(t).
\end{equation}
元のゼータ関数に $(1-2^{-t})$ と $(1-3^{-t})$ をかけることで，
素数 $2$ と $3$ のオイラー因子を打ち消しています．

\subsection{なぜ重力が消えるのか}

ここが定理のキモです．$t = 0$ のとき:
\begin{align}
  1 - 2^{-t}\Big|_{t=0} &= 1 - 2^0 = 1 - 1 = 0, \\
  1 - 3^{-t}\Big|_{t=0} &= 1 - 3^0 = 1 - 1 = 0.
\end{align}
どちらも\textbf{ゼロ}！ 一方，$\zeta(0) = -1/2$ は
ゼロではない有限の値です．

したがって:
\[
  K_{\neg\{2,3\}}(t) \approx (t \log 2)(t \log 3) \times (-\tfrac{1}{2})
  = -\frac{\log 2 \cdot \log 3}{2}\, t^2 \quad (t \to 0).
\]
$t^2$ に比例してゼロに近づく！

\begin{tcolorbox}[colback=red!5,colframe=red!50!black,title=定理1: 重力無効化]
2つ以上の素数 $\{p_1, p_2, \ldots\}$ をミュートすると，
修正ヒートカーネルは $t = 0$ で\textbf{2重ゼロ}を持つ．

その結果:
\begin{itemize}
  \item 宇宙項（宇宙の膨張エネルギー）: $a_0 = 0$（消滅）
  \item アインシュタインの重力項: $a_2 = 0$（消滅）
  \item \textbf{有効重力定数: $G_{\text{eff}} = 0$}
\end{itemize}
つまり，\textbf{重力が完全にゼロ}になる！
\end{tcolorbox}

\subsection{ノイズキャンセリング・ヘッドフォンのアナロジー}

これは\textbf{ノイズキャンセリング・ヘッドフォン}と似ています．

ヘッドフォンのマイクが外部の騒音を拾い，
逆位相の音波を出すことで音が打ち消されますよね．
2つの波がぴったり重なって，結果は「沈黙」になる．

同じように:
\begin{itemize}
  \item 素数 $2$ のフィルター: 重力の一部を打ち消す波を生成
  \item 素数 $3$ のフィルター: 残りを打ち消す波を生成
  \item 2つ合わせると: \textbf{重力の源が完全に消える}
\end{itemize}

\begin{tcolorbox}[colback=gray!5,colframe=gray!50!black]
\textbf{直感的まとめ}: 重力は無限個の素数チャンネルの
「合奏」として生まれている．
2つのチャンネルにフィルターをかけると，
$t = 0$（紫外端）で合奏全体が沈黙する．
結果: $G_{\text{eff}} = 0$．
\end{tcolorbox}

%% ============================================================
\section{定理2: 反重力 --- 全素数を$\pi$回転すると$G_{\text{eff}} = -2G$}
%% ============================================================

\subsection{$\pi$回転とは}

今度は素数を消すのではなく，
すべての素数チャンネルに「位相を $\pi$（180度）回転」
させる操作を考えます．

オイラー因子の変換:
\[
  \frac{1}{1 - p^{-t}} \quad\longrightarrow\quad \frac{1}{1 + p^{-t}}.
\]
符号が反転しただけ！ しかしその効果は劇的です．

全素数で $\pi$ 回転した新しいヒートカーネルは:
\begin{equation}
  K_\pi(t) = \prod_{p} \frac{1}{1 + p^{-t}} = \frac{\zeta(2t)}{\zeta(t)}.
\end{equation}

\subsection{リウヴィル関数: すべてのモードに ``フェルミオン符号''}

この変換は，数論の\textbf{リウヴィル関数} $\lambda(n)$ と
深い関係があります．

$\lambda(n) = (-1)^{\Omega(n)}$，
ここで $\Omega(n)$ は $n$ の素因数の個数（重複込み）です．

\begin{center}
\begin{tabular}{c|cccccccc}
$n$ & 1 & 2 & 3 & 4 & 5 & 6 & 7 & 8 \\
\hline
$\Omega(n)$ & 0 & 1 & 1 & 2 & 1 & 2 & 1 & 3 \\
$\lambda(n)$ & $+1$ & $-1$ & $-1$ & $+1$ & $-1$ & $+1$ & $-1$ & $-1$ \\
\end{tabular}
\end{center}

素因数が奇数個なら $-1$（マイナス），偶数個なら $+1$（プラス）．
物理学では，こういう「$\pm 1$ の符号」を持つ粒子を
\textbf{フェルミオン}（電子のような粒子）と呼びます．
つまり $\pi$ 回転は，
\textbf{真空のすべてのモードにフェルミオン的な符号をつける}操作です．

\subsection{重力が逆転する}

計算すると:
\begin{align}
  K_\pi(0) &= \frac{\zeta(0)}{\zeta(0)} = 1, \\
  K_\pi'(0) &= \log(2\pi).
\end{align}
一方，通常のヒートカーネルでは
$\zeta'(0) = -\frac{1}{2}\log(2\pi)$ です．

重力定数の比:
\begin{equation}
  \frac{G_{\text{eff}}}{G}
  = \frac{K_\pi'(0)}{\zeta'(0)}
  = \frac{\log(2\pi)}{-\frac{1}{2}\log(2\pi)}
  = -2.
\end{equation}

\begin{tcolorbox}[colback=purple!5,colframe=purple!50!black,title=定理2: 重力逆転]
すべての素数チャンネルに $\pi$ 回転（位相反転）を適用すると:
\[
  G_{\text{eff}} = -2G.
\]
重力の符号が反転し，強さは2倍になる．
つまり，\textbf{物体は引き合うのではなく，押し合う}（反重力）！
\end{tcolorbox}

\subsection{音楽を逆再生するアナロジー}

これは\textbf{音楽を逆再生する}のに似ています．

録音された曲のすべての音の波形を上下反転（位相反転）すると，
曲の「雰囲気」が変わります．
メジャーコード（明るい響き）がマイナーコード（暗い響き）に
聞こえるような変化です．

同様に:
\begin{itemize}
  \item 通常の真空: すべてのモードが ``ボソン的''（引力を生む）
  \item $\pi$ 回転後: すべてのモードが ``フェルミオン的''（斥力を生む）
  \item 結果: 重力の方向が逆転する
\end{itemize}

%% ============================================================
\section{ワープバブル: 外殻（反重力）+ 内部（無重力）}
%% ============================================================

2つの定理を組み合わせると，
SF映画のような「ワープバブル」が構築できます！

\subsection{二層構造}

\begin{center}
\begin{tabular}{lccc}
\toprule
層 & 操作 & $G_{\text{eff}}$ & 光速 \\
\midrule
外殻（シェル） & 全素数を$\pi$回転 & $-2G$（反重力） & $c$（$d_S \approx 2$） \\
内部（船室） & 素数 $\{2,3\}$ をミュート & $0$（無重力） & $c$（$d_S = 4$） \\
外部（宇宙） & 操作なし & $G$（通常重力） & $c$（$d_S = 4$） \\
\bottomrule
\end{tabular}
\end{center}

\begin{tcolorbox}[colback=cyan!5,colframe=cyan!50!black,title=アスキーアート: ワープバブルの断面図]
\begin{verbatim}
      外部: 通常の宇宙  G = G
   ┌──────────────────────────┐
   │  外殻: 反重力シェル       │
   │  G_eff = -2G             │
   │  (全素数π回転)            │
   │  d_S ≈ 2 (ホログラフィック) │
   │  ┌──────────────────┐    │
   │  │  内部: 無重力空間  │    │
   │  │  G_eff = 0        │    │
   │  │  (素数2,3ミュート)  │    │
   │  │  c_eff = c (光速)  │    │
   │  │  d_S = 4 (通常)    │    │
   │  │                    │    │
   │  │    🚀 宇宙船       │    │
   │  │                    │    │
   │  └──────────────────┘    │
   └──────────────────────────┘
\end{verbatim}
\end{tcolorbox}

\subsection{内部では光速は変わらない}

定理1で重力をゼロにした内部空間では，
\textbf{光の速さは通常どおり $c$}（秒速30万km）です．

これは「コンヌの距離公式」という定理で証明されます．
素数のミュートは「ポテンシャル $V$」の追加に対応しますが，
距離の計算には影響しない（$[V, f] = 0$ という性質のため）
からです．

宇宙船の中にいる人は，普通の物理法則のもとで
普通に生活できます．

\subsection{外殻のホログラフィック性}

反重力シェルでは面白いことが起きます．
通常は4次元（3次元空間 + 1次元時間）の時空が，
\textbf{実効的に2次元に縮退}します（$d_S \approx 2$）．

これは「ホログラフィック次元縮小」と呼ばれ，
シェルの壁が2次元の面のように振る舞うことを意味します．
2次元の因果構造は4次元の光速に縛られないため，
これが超光速的な伝搬メカニズムを提供する可能性があります．

%% ============================================================
\section{本当にできるの？ 実験の話}
%% ============================================================

「数学的にはわかった．でも本当に物理的にできるの？」
---当然の疑問です．

この理論には，\textbf{たった1つの実験で白黒がつく}
鋭い予測があります．

\subsection{BAW結晶実験}

\begin{tcolorbox}[colback=yellow!5,colframe=yellow!60!black,title=実験の概要]
\textbf{装置}: 142\,$\mu$m の水晶BAW（体積弾性波）結晶

\textbf{やること}: 素数 $2$ と $3$ のモードだけを
フィルターで除いた音響キャビティを作り，
カシミール・エネルギー（真空のエネルギー）の比を測定する

\textbf{費用}: 約33万円（\$2{,}200）

\textbf{期間}: 4--6ヶ月
\end{tcolorbox}

\subsection{182 vs 1/3 テスト}

この実験の予測値は:
\begin{equation}
  \frac{|E_{\neg\{2,3\}}|}{|E_{\text{通常}}|}
  = |(1 - 2^3)(1 - 3^3)| = |{-7}| \times |{-26}| = 182.
\end{equation}

一方，もし素数の構造が物理に関係ないなら（オイラー積ではなく
単純なモード数え上げ），予測値は $1/3$ になります．

\begin{tcolorbox}[colback=green!5,colframe=green!50!black,title=判定基準]
\begin{itemize}
  \item 測定値 $\approx 182$（誤差50\%以内）:
    \textbf{数論は物理学だ！}
    定理1と定理2は物理的に実現可能．
    ワープバブルの基盤が実験的に確認される．
  \item 測定値 $\approx 1/3$:
    オイラー積は数学的には美しいが，
    物理的真空には直接対応しない．
\end{itemize}
2つの予測値の比は $182 \div (1/3) = 546$ 倍．
\textbf{曖昧さの余地はほぼない}．
\end{tcolorbox}

%% ============================================================
\section{正直な評価: 何がわかっていて何がわかっていないか}
%% ============================================================

最後に，この研究の現状を正直に整理します．

\subsection{わかっていること（数学的に証明済み）}

\begin{itemize}
  \item \textbf{定理1は数学的に厳密に正しい}:
    2つの素数をミュートすると，ヒートカーネルは $t=0$ で
    2重ゼロを持ち，$a_0 = a_2 = 0$ となる．
  \item \textbf{定理2も数学的に厳密に正しい}:
    全素数の $\pi$ 回転で $G_{\text{eff}}/G = -2$ になる．
  \item \textbf{光速不変}: $G=0$ 領域で $c_{\text{eff}} = c$
    はコンヌの距離公式から証明される．
  \item \textbf{ホログラフィック縮退}: 反重力シェルの
    スペクトル次元 $d_S \approx 2$ は計算で確認済み．
\end{itemize}

\subsection{わかっていないこと（未証明）}

\begin{itemize}
  \item \textbf{物理的実現可能性は未証明}:
    Bost-Connes系のスペクトル三重が物理的真空を
    忠実に記述するかどうかは，まだ仮説の段階．
  \item \textbf{空間構造の問題}:
    Bost-Connes系は $(0+1)$ 次元（時間だけ）のシステム．
    3+1次元の時空への埋め込みは完全には定式化されていない．
  \item \textbf{非摂動論的補正}:
    $t = 0$ 付近の展開に，まだ捉えきれていない補正が
    あるかもしれない．
\end{itemize}

\begin{tcolorbox}[colback=orange!5,colframe=orange!60!black,title=一番大事なこと]
\textbf{定理は数学的に厳密．しかし物理的に実現できるかは未証明．}

そして，それを決めるのは\textbf{実験}です．

33万円のBAW結晶実験が $182$ を出すか $1/3$ を出すか
--- それだけで，この理論の運命が決まります．

もし $182$ なら:
素数の数学的構造が，真空のエネルギーを直接支配していることになり，
重力制御への道が開かれます．

もし $1/3$ なら:
美しい数学ではあるが，自然はそのようには動かない，
ということになります．

\textbf{すべては実験が決める}．
これこそが科学の本質です．
\end{tcolorbox}

\vspace{1em}

\begin{center}
\textit{--- 数学は言語であり，実験は判事である ---}
\end{center}

\end{document}
